\documentclass[aspectratio=169,11pt]{beamer}
\usepackage[T1]{fontenc}
\usepackage[utf8]{inputenc}
\usepackage[provide=*,polish]{babel}
\usepackage{lmodern}
\usepackage{amsmath}
\usepackage{booktabs}
\usepackage{graphicx}
\usepackage{hyperref}

\definecolor{navy}{HTML}{152A43}
\definecolor{teal}{HTML}{007F86}
\definecolor{light}{HTML}{EFF5F6}
\setbeamercolor{normal text}{fg=navy,bg=white}
\setbeamercolor{frametitle}{fg=navy,bg=white}
\setbeamercolor{structure}{fg=teal}
\setbeamercolor{block title}{fg=white,bg=teal}
\setbeamercolor{block body}{fg=navy,bg=light}
\setbeamertemplate{navigation symbols}{}
\setbeamertemplate{footline}[frame number]
\setbeamertemplate{itemize items}[circle]
\setbeamersize{text margin left=11mm,text margin right=11mm}

\title{MuJoCo HalfCheetah}
\subtitle{Uczenie ze wzmocnieniem \quad | \quad Pierwszy punkt kontrolny}
\author{Aleksandra Wieczorek \and Wojciech W\'ojcik  }
\institute{Politechnika Warszawska}
\date{29 pa\'zdziernika 2026}

\begin{document}

\begin{frame}
  \titlepage
  \vspace{-0.5em}
  \begin{center}\small Temat zielony\end{center}
\end{frame}

\begin{frame}{Cel zadania}
  \begin{columns}[T,onlytextwidth]
    \begin{column}{0.61\textwidth}
      \begin{itemize}
        \item Sterowanie dwuwymiarowym modelem geparda w symulatorze fizyki MuJoCo.
        \item Agent wybiera momenty steruj\k{a}ce sze\'scioma stawami.
        \item Zadanie: porusza\'c si\k{e} mo\.zliwie szybko w prawo, ograniczaj\k{a}c koszt sterowania.
        \item Wariant: Gymnasium \texttt{HalfCheetah-v5}.
      \end{itemize}
    \end{column}
    \begin{column}{0.35\textwidth}
      \begin{block}{Charakter problemu}
        1 agent\\[0.3em]
        Ci\k{a}g\l e akcje\\[0.3em]
        Symulacja fizyczna\\[0.3em]
        Brak przeciwnik\'ow
      \end{block}
    \end{column}
  \end{columns}
\end{frame}

\begin{frame}{P\k{e}tla interakcji ze \'srodowiskiem}
  \centering
  \begin{block}{Jeden krok symulacji}
    \centering
    Obserwacja $s_t$ $\longrightarrow$ polityka $\pi$ $\longrightarrow$ akcja $a_t$\\[0.6em]
    \centering $\longrightarrow$ MuJoCo $\longrightarrow$ $s_{t+1}$ oraz $r_t$
  \end{block}
  \vspace{1em}
  \begin{itemize}
    \item Domy\'slny epizod: \textbf{1000 krok\'ow}; ograniczenie czasowe przerywa epizod.
  \end{itemize}
\end{frame}

\begin{frame}{Co obserwuje agent?}
  \begin{columns}[T,onlytextwidth]
    \begin{column}{0.34\textwidth}
      \begin{block}{Przestrze\'n obserwacji}
        \centering $s_t\in\mathbb{R}^{17}$
      \end{block}
    \end{column}
    \begin{column}{0.62\textwidth}
        \begin{itemize}
        \item 8 warto\'sci po\l o\.ze\'n i k\k{a}t\'ow (\texttt{qpos}).
        \item 9 pr\k{e}dko\'sci (\texttt{qvel}).
        \item Opcjonalnie dodatkowy parametr położenia x.
        \item Obserwacja zawiera dane stanu fizycznego, a nie obraz z kamery.
      \end{itemize}
    \end{column}
  \end{columns}
\end{frame}

\begin{frame}{Jakie akcje mo\.ze wykona\'c agent?}
  \begin{block}{Ci\k{a}g\l a przestrze\'n sterowania}
    \centering $a_t\in[-1,1]^6$
  \end{block}
  \begin{columns}[T,onlytextwidth]
    \begin{column}{0.48\textwidth}
      \textbf{Tylne ko\'nczyny}
      \begin{enumerate}
        \item Udo (\texttt{bthigh})
        \item Podudzie (\texttt{bshin})
        \item Stopa (\texttt{bfoot})
      \end{enumerate}
    \end{column}
    \begin{column}{0.48\textwidth}
      \textbf{Przednie ko\'nczyny}
      \begin{enumerate}\setcounter{enumi}{3}
        \item Udo (\texttt{fthigh})
        \item Podudzie (\texttt{fshin})
        \item Stopa (\texttt{ffoot})
      \end{enumerate}
    \end{column}
  \end{columns}
  \vspace{0.4em}
  \small Ka\.zda sk\l adowa okre\'sla sterowanie momentem si\l y w stawie. Warto\'sci s\k{a} ograniczone do $[-1,1]$.
\end{frame}

\begin{frame}{Funkcja nagrody}
  \begin{block}{Nagroda za pojedynczy krok}
    \centering \Large $r_t = w_f\frac{x_{t+1}-x_t}{\Delta t}-w_c\lVert a_t\rVert_2^2$
  \end{block}
  \vspace{0.3em}
  \begin{itemize}
    \item $w_f=1$: dodatnia nagroda za pr\k{e}dko\'s\'c do przodu.
    \item $w_c=0{,}1$: kara za du\.ze sygna\l y steruj\k{a}ce.
    \item Domy\'slnie $\Delta t=5\cdot0{,}01=0{,}05$ s.
    \item W eksperymentach rejestrujemy osobno oba sk\l adniki.
  \end{itemize}
\end{frame}

\begin{frame}{Mo\.zliwo\'sci konfiguracji}
  \small
  \begin{tabular}{@{}p{0.47\textwidth}p{0.18\textwidth}p{0.26\textwidth}@{}}
    \toprule
    Parametr & Domy\'slnie & Znaczenie\\
    \midrule
    \texttt{forward\_reward\_weight} & 1,0 & Waga pr\k{e}dko\'sci\\
    \texttt{ctrl\_cost\_weight} & 0,1 & Kara za akcje\\
    \texttt{reset\_noise\_scale} & 0,1 & Losowo\'s\'c startu\\
    \texttt{frame\_skip} & 5 & Odst\k{e}p sterowania\\
    \texttt{exclude\_current\_positions\_}\\\texttt{from\_observation} & true & Dost\k{e}p do pozycji $x$\\
    \texttt{max\_episode\_steps} & 1000 & Limit epizodu\\
    \bottomrule
  \end{tabular}
  \vspace{0.6em}
\end{frame}

\begin{frame}{Punkt odniesienia i metryki}
  \begin{columns}[T,onlytextwidth]
    \begin{column}{0.45\textwidth}
      \begin{block}{Punkty odniesienia}
        \begin{itemize}
          \item Teraz: agent losowy.
          \item Grudzie\'n: SAC jako podstawowy agent uczony.
          \item Opcjonalnie: PPO.
        \end{itemize}
      \end{block}
    \end{column}
    \begin{column}{0.51\textwidth}
      \begin{block}{Metryki ewaluacji}
        \begin{itemize}
          \item Zwrot z epizodu.
          \item Pr\k{e}dko\'s\'c i przemieszczenie $x$.
          \item $\frac{1}{T}\sum_t\lVert a_t\rVert_2^2$.
          \item Zmienno\'s\'c wynik\'ow mi\k{e}dzy ziarnami.
        \end{itemize}
      \end{block}
    \end{column}
  \end{columns}
  \vspace{0.2em}
  \small Dla zmienionej funkcji nagrody por\'ownujemy r\'ownie\.z metryki niezale\.zne od jej wag.
\end{frame}

\begin{frame}{Plan dalszych prac}
  \begin{enumerate}
    \setlength{\itemsep}{0.6em}
    \item Trening SAC, odtwarzalne uruchomienia, wst\k{e}pne por\'ownanie z agentem losowym.
    \item Eksperymenty kontrolowane, kilka wariant\'ow parametr\'ow i ziaren losowo\'sci.
    \item Raport, analiza i prezentacja wynik\'ow.
  \end{enumerate}
\end{frame}

\end{document}
